\section{Limitations}
\label{sec:limitations}

NTD-PL assumes that a common nonlinear response explains a coherent
part of the residual left by matched-rank Tucker. This scalar link may be
less effective when nonlinearities vary strongly across spatial regions,
spectral bands, modes, or samples. In such cases, a single polynomial response
can underfit the domain variation even if the Tucker backbone is appropriate.

NTD-PL enhances a fixed Tucker rank budget by changing the measurement map on
top of the latent tensor. When the remaining error is ordinary multilinear
structure rather than a scalar response, increasing Tucker rank remains the
more direct modeling choice.

Polynomial degree introduces a second tradeoff. Higher degrees increase the
approximation capacity of the link, but they can amplify Tucker scale
ambiguity and numerical conditioning issues in the power basis. This motivates
the ridge-regularized Link Refresh step, factor normalization, and degree
continuation used in the optimizer.

Finally, the joint problem remains nonconvex. The Link Refresh subproblem is a
small ridge regression once the latent tensor is fixed, but the Tucker backbone
and polynomial link are learned jointly through a nonconvex objective. We do not
claim global optimality.

Our strongest empirical evidence is currently from controlled nonlinear
residuals and hyperspectral reconstruction and completion. Broader validation
on large sparse tensors, temporal or multimodal tensors, and domain-varying
nonlinear responses is left for future work.
